\newpage

% *************************
% **     PROBLEMA        **
% *************************

\section{Problema de Investigación}
\subsection{Descripción del problema}La educación rural bajo el modelo de Educación Intercultural Bilingüe (EIB) en el Perú enfrenta una profunda brecha de recursos y herramientas tecnológicas adaptadas. A pesar de que el sistema EIB atiende a más de 1,2 millones de estudiantes en el país \citep{minedu2018}, su implementación en el ámbito rural se ve limitada por la falta de conectividad digital: según el INEI, apenas el 20,5\% de los hogares rurales tiene acceso a internet \citep{inei2025}. Esta desconexión digital afecta de manera crítica a los departamentos del sur andino (Cusco y Puno), donde se concentra una alta proporción de población escolar quechuahablante.

En este contexto regional, la lengua originaria predominante es el Quechua Collao. Esta variedad lingüística impone desafíos complejos para los sistemas de procesamiento automático: su naturaleza aglutinante y sufijadora genera palabras de alta densidad semántica, y en la práctica educativa, los estudiantes y docentes interactúan de forma dinámica alternando constantemente entre el quechua y el español (translenguaje o alternancia de códigos). Las herramientas tecnológicas estándar no están preparadas para indexar, segmentar o tokenizar esta variedad lingüística de bajos recursos digitales, ni para discernir entre las consultas formuladas en una u otra lengua.

Por otro lado, aunque el avance de la inteligencia artificial generativa ha popularizado el uso de asistentes virtuales para el aprendizaje, las soluciones comerciales son inviables para las escuelas rurales EIB. Esto se debe a tres limitaciones estructurales: (a) requieren conexión permanente a internet para consultar modelos remotos en la nube, (b) carecen de representación real y alineación con la variedad Quechua Collao, y (c) no están adaptadas al currículo oficial del Ministerio de Educación ni siguen estrategias instruccionales de andamiaje formativo.

En consecuencia, existe una brecha técnica y pedagógica no resuelta: la ausencia de un asistente educativo bilingüe español--Quechua Collao que funcione 100\% offline, que integre un enrutamiento lingüístico automático para consultas bilingües, y que recupere localmente contenidos curriculares validados mediante una arquitectura de modelo compacto y RAG local. Esta brecha justifica el desarrollo y evaluación del prototipo propuesto.

\subsection{Identificación del problema}
El problema central de la investigación es la ausencia de un asistente educativo bilingüe español--Quechua Collao que funcione sin conexión a internet, basado en un modelo compacto con RAG local, para responder consultas pedagógicas en contextos de Educación Intercultural Bilingüe.

\subsection{Formulación del problema}
\subsubsection{Problema general}
¿Cómo diseñar, implementar y evaluar un asistente educativo offline español--Quechua Collao basado en un modelo compacto con RAG local para apoyar consultas curriculares en contextos de Educación Intercultural Bilingüe?

\subsubsection{Subproblemas}
\begin{itemize}
	\item ¿Cuáles son los requerimientos lingüísticos, pedagógicos y tecnológicos que debe satisfacer un asistente educativo bilingüe offline para contextos rurales de EIB?
	\item ¿Cómo diseñar una arquitectura offline que integre detección automática del idioma, recuperación local de contenidos curriculares, recursos lingüísticos explícitos y generación de respuestas pedagógicas?
	\item ¿Cómo implementar y configurar un modelo compacto cuantizado para operar de forma eficiente sin conexión a internet en dispositivos de bajo costo?
	\item ¿Qué diferencia en rendimiento y pertinencia se registra al comparar tres configuraciones básicas del sistema: modelo base sin RAG, modelo con RAG local y modelo con el enfoque híbrido completo?
	\item ¿Cómo evaluar la calidad del prototipo resultante mediante métricas técnicas, lingüísticas y rúbricas de pertinencia pedagógica?
\end{itemize}
\section{Justificación}
La investigación se justifica porque las condiciones descritas en el problema —escasa conectividad rural, ausencia de herramientas tecnológicas en Quechua Collao y más de 1,2 millones de estudiantes EIB desatendidos tecnológicamente \citep{minedu2018}— no se resuelven suficientemente con las herramientas de inteligencia artificial generativa comerciales en la nube, las cuales dependen de conectividad permanente y operan mayoritariamente sobre lenguas de alta disponibilidad. El proyecto propone una alternativa a esa brecha desde tres perspectivas:

\begin{itemize}
	\item \textbf{Justificación Teórica:} El estudio aporta a la integración del procesamiento de lenguas de bajos recursos digitales con la tecnología educativa. Aunque existen investigaciones previas sobre traducción automática o evaluación de modelos para quechua de manera aislada, no se ha propuesto un diseño que acople el entendimiento bilingüe contextualizado, la recuperación local de contenidos y mecanismos de adaptación orientados al andamiaje pedagógico para el Quechua Collao.
	\item \textbf{Justificación Práctica:} Propone un prototipo funcional offline de apoyo pedagógico complementario en áreas rurales de baja o nula conectividad, facilitando el acceso a materiales oficiales en la lengua originaria del estudiante.
	\item \textbf{Justificación Metodológica:} Diseña un marco de evaluación comparativo que permite medir el aporte real de cada componente del sistema (el modelo base, la recuperación local de contenidos y la alineación de estilo del modelo), combinando métricas automatizadas de rendimiento y calidad de texto con valoraciones cualitativas aplicadas por especialistas bilingües.
\end{itemize}

% *************************
% **     ALCANCES        **
% *************************
\section{Alcances y limitaciones}
\subsection{Alcances}
El presente estudio abarca el diseño, implementación y evaluación de un prototipo de asistente educativo bilingüe español--Quechua Collao basado en un modelo compacto de lenguaje. La investigación se centra en la asistencia textual para estudiantes y docentes de contextos de EIB, considerando los siguientes alcances específicos:

\begin{itemize}
	\item \textbf{Alcance lingüístico:} Se delimita de forma exclusiva al Quechua Collao, de acuerdo con la ubicación geográfica del estudio y la disponibilidad de corpus paralelos como los construidos por \citet{Orcotoma2025}. Esta variedad presenta rasgos fonológicos diferenciadores (consonantes oclusivas glotizadas y aspiradas) que deben reflejarse cuidadosamente en los recursos del sistema. Esta delimitación evita prometer cobertura de otras variedades dialectales, permitiendo construir un prototipo de alta pertinencia regional.

	\item \textbf{Alcance pedagógico:} El primer ciclo de validación se centrará en consultas curriculares de 3.° a 6.° grado de primaria, priorizando las áreas de Comunicación y Ciencia y Tecnología según la disponibilidad de materiales EIB bilingües. Las respuestas del sistema seguirán el principio de andamiaje cognitivo, priorizando explicaciones graduales sobre respuestas directas.

	\item \textbf{Alcance tecnológico:} Comprende una arquitectura offline compuesta por detección automática del idioma, enrutamiento lingüístico, generación aumentada por recuperación (RAG) local y un modelo compacto de lenguaje cuantizado, preferentemente a 4 bits según las pruebas de rendimiento realizadas. El aporte tecnológico radica en la integración, configuración y evaluación comparativa de esta arquitectura.

	\item \textbf{Alcance disciplinar:} Se ubica en la convergencia de la ingeniería informática, el procesamiento de lenguaje natural y la tecnología educativa para contextos bilingües interculturales.
\end{itemize}

\subsection{Limitaciones}
Se reconocen las siguientes limitaciones que condicionan la ejecución del proyecto:

\begin{itemize}
	\item \textbf{Disponibilidad de datos:} Al tratarse de una lengua de bajos recursos digitales, los datos bilingües se generarán con modelos de IA y se validarán contra corpus de referencia existentes \citep{Orcotoma2025, Duran2026}. La calidad de estos datos es crítica para el funcionamiento del sistema.
	\item \textbf{Cuello de botella de la tokenización:} La morfología aglutinante del quechua genera una fertilidad de tokens estimada en $\approx$3,4$\times$ respecto al español, lo que penaliza la latencia en CPU e infla el consumo de la caché KV del modelo con contextos extensos.
	\item \textbf{Restricción dialectal:} El prototipo está orientado exclusivamente al Quechua Collao, por lo que el sistema no está diseñado para generalizarse a otras variedades dialectales.
	\item \textbf{Restricciones de hardware:} La ejecución offline de un modelo compacto depende de la memoria disponible (RAM) del dispositivo local, limitando el tamaño del modelo utilizable y la latencia aceptable de generación \citep{Xu2024}.
	\item \textbf{Validación pedagógica y lingüística:} La evaluación de pertinencia y adecuación dependerá de la validación del docente especialista EIB y del especialista lingüístico en la variedad Quechua Collao.
	\item \textbf{Alcance textual:} El prototipo no incorpora reconocimiento ni síntesis de voz, considerando la asistencia por audio fuera del alcance de este trabajo.
\end{itemize}

% *************************
% **     OBJETIVOS       **
% *************************

\section{Objetivos}
\subsection{Objetivo General}
Diseñar, implementar y evaluar un prototipo de asistente educativo offline español--Quechua Collao basado en un modelo de lenguaje compacto cuantizado, generación aumentada por recuperación local y recursos lingüísticos validados, orientado al apoyo de consultas curriculares en contextos rurales de Educación Intercultural Bilingüe.

\subsection{Objetivos Específicos}
\begin{itemize}
	\item \textbf{Curar y estructurar} una base documental bilingüe español--Quechua Collao, integrando materiales curriculares oficiales, glosarios terminológicos, diccionarios locales y datos generados con IA validados contra corpus de referencia.
	\item \textbf{Diseñar} una arquitectura offline que integre detección automática del idioma, recuperación local de información y generación de respuestas pedagógicas.
	\item \textbf{Implementar} un prototipo funcional del asistente educativo seleccionando, cuantizando y configurando el modelo compacto base y el pipeline de recuperación.
	\item \textbf{Comparar} configuraciones básicas del sistema para medir el aporte real de la recuperación local, las reglas de normalización y la adaptación de estilo del modelo.
	\item \textbf{Evaluar} el prototipo mediante métricas técnicas de rendimiento offline, fidelidad lingüística mediante métricas automáticas y pertinencia pedagógica a través de rúbricas de juicio de expertos.
	\item \textbf{Documentar} las limitaciones de la arquitectura propuesta, los riesgos de transferencia dialectal y las recomendaciones de mejora.
\end{itemize}

\subsection{Alineación con los Objetivos de Desarrollo Sostenible}
La investigación se alinea con los siguientes Objetivos de Desarrollo Sostenible (ODS):

\begin{itemize}
	\item \textbf{ODS 4 — Educación de calidad.} El estudio contribuye a una educación inclusiva y equitativa al proveer una herramienta de apoyo pedagógico en lengua materna para estudiantes de zonas rurales bajo el modelo EIB.
	\item \textbf{ODS 9 — Industria, innovación e infraestructura.} Promueve la innovación tecnológica al diseñar una solución de inteligencia artificial offline, adaptada a condiciones de hardware de bajo costo e infraestructura digital limitada.
	\item \textbf{ODS 10 — Reducción de las desigualdades.} Aborda la brecha digital y lingüística que afecta a comunidades quechuahablantes, facilitando el acceso a recursos educativos interactivos sin requerir conexión a internet.
\end{itemize}
